% =============================================================
%  3. PRUEBAS EN EL PASILLO (común)
% =============================================================
\renewcommand{\seccorta}{Pruebas}
\section{Pruebas en el pasillo de la facultad}

\begin{frame}{Ubicación del enlace y línea base}
  \begin{columns}[c]
    \column{0.42\textwidth}\centering
    \figura{plano_pasillo.jpg}{\linewidth}\\
    {\scriptsize Enlace de 70 m entre A y B}
    \column{0.55\textwidth}
    \small
    \begin{itemize}
      \item El laboratorio era demasiado chico para un equipo pensado para varios km:
            se pasó al pasillo (\textbf{70 m}).
      \item Misma configuración (5540 MHz, 60 MHz, roles).
      \item Sin análisis de espectro en el pasillo: con el Punto B gestionado por Wi-Fi,
            airView queda deshabilitado.
      \item Línea base: $-29$/$-30$ dBm, $\Delta$ entre cadenas de 2 a 4 dB.
    \end{itemize}
  \end{columns}
\end{frame}

\begin{frame}{Comportamiento del enlace en el pasillo}
  \begin{columns}[c]
    \column{0.5\textwidth}\centering
    \figura{panel_pasillo.jpg}{\linewidth}
    \column{0.47\textwidth}
    \small
    \begin{block}{Capacidad y tráfico}
      Capacidad de 443--454 Mbps (256QAM, 8x). Latencia base de 1 ms con picos de hasta
      2.4 ms por \emph{multipath} en las paredes.
    \end{block}
    \begin{block}{Señal y ruido}
      Caídas abruptas de $-28$ a $-56$ dBm por el paso de personas que obstruyen la primera
      zona de Fresnel. Piso de ruido estable en $-86$ dBm.
    \end{block}
  \end{columns}
\end{frame}

\begin{frame}{Alineación y ajuste fino}
  \begin{columns}[c]
    \column{0.22\textwidth}\centering
    \figura{alineadas.jpg}{\linewidth}\\ {\scriptsize Alineadas}
    \column{0.26\textwidth}\centering
    \figura{constelacion_alineadas.jpg}{\linewidth}\\ {\scriptsize 256QAM}
    \column{0.22\textwidth}\centering
    \figura{desalineadas.jpg}{\linewidth}\\ {\scriptsize Desalineadas ($\pm$90°)}
    \column{0.26\textwidth}\centering
    \figura{constelacion_desalineadas.jpg}{\linewidth}\\ {\scriptsize QPSK}
  \end{columns}
  \vspace{4pt}
  \centering\small
  \begin{tabular}{@{}lcc@{}}
    \toprule
     & \textbf{Alineadas} & \textbf{Desalineadas} \\ \midrule
    Señal local / remota & $-36$ / $-35$ dBm & $-54$ / $-52$ dBm \\
    $\Delta$ entre cadenas (local / remoto) & 8 / 5 dB & 6 / 3 dB \\
    Modulación & 256QAM & QPSK \\
    \bottomrule
  \end{tabular}
\end{frame}

\begin{frame}{Interpretación de la alineación}
  \begin{itemize}
    \item \textbf{Alineadas:} mucha potencia y 256QAM, pero con $\Delta=8$ dB: el
          \emph{multipath} del piso y las paredes afecta distinto a cada polarización.
          La constelación dispersa muestra el aumento del EVM.
    \item \textbf{Desalineadas:} la señal cae $\approx$18 dB (lóbulos secundarios y rebotes) y
          el algoritmo adaptativo baja a \textbf{QPSK}, priorizando robustez sobre capacidad.
          El $\Delta$ baja porque ambas polarizaciones se degradan por igual.
    \item \textbf{En ambos casos:} lecturas fluctuantes por el tránsito de estudiantes que
          obstruye la primera zona de Fresnel.
  \end{itemize}
\end{frame}

\begin{frame}{Pruebas de latencia (5 paquetes ICMP de 56 bytes)}
  \begin{columns}[c]
    \column{0.48\textwidth}\centering
    \figura{ping_alineadas.jpg}{0.75\linewidth}\\ {\scriptsize Alineadas}
    \column{0.48\textwidth}\centering
    \figura{ping_desalineadas.jpg}{0.75\linewidth}\\ {\scriptsize Desalineadas}
  \end{columns}
  \vspace{4pt}
  \centering\small
  \begin{tabular}{@{}lccc@{}}
    \toprule
     & \textbf{Mín.} & \textbf{Prom.} & \textbf{Máx.} \\ \midrule
    Alineadas (256QAM) & 1.133 ms & 2.93 ms & \textcolor{red}{8.103 ms} \\
    Desalineadas (QPSK) & 2.103 ms & 3.13 ms & 5.45 ms \\
    \bottomrule
  \end{tabular}\\[3pt]
  \raggedright\small 256QAM reduce el tiempo de aire (menor mínimo), pero es más sensible a
  obstrucciones y rebotes (mayor jitter). QPSK sube la latencia base y amortigua los picos.
\end{frame}

\begin{frame}{Prueba de velocidad (Speed Test integrado)}
  \begin{columns}[c]
    \column{0.45\textwidth}\centering
    \figura{speedtest.jpg}{\linewidth}
    \column{0.52\textwidth}
    \small
    \begin{block}{Resultado}
      RX máx.\ \textbf{266.96 Mbps} y ping de 1.35 ms, frente a 454 Mbps de capacidad
      ($\approx$40\,\% menos).
    \end{block}
    \begin{itemize}
      \item \textbf{CPU de la antena:} el Speed Test genera el tráfico en el propio equipo y
            el procesador se satura antes que el radio.
      \item \textbf{Overhead:} la capacidad es bruta; el test mide la tasa útil
            (encabezados IP/MAC y cifrado WPA2).
    \end{itemize}
    $\Rightarrow$ La limitación es de la herramienta, no del enlace de RF.
  \end{columns}
\end{frame}

\begin{frame}{Pruebas de estrés con iPerf3}
  \begin{alertblock}{No se pudo realizar}
    iPerf3 necesita una PC cableada en cada extremo (cliente--servidor), y el Punto B se
    gestionaba solo desde el celular.
  \end{alertblock}
  \vspace{4pt}
  \begin{columns}[T]
    \column{0.48\textwidth}
    \begin{block}{TCP}
      \small Con acuses de recibo (ACK): mide la velocidad útil real, incluyendo
      retransmisiones por un entorno inestable.
    \end{block}
    \column{0.48\textwidth}
    \begin{block}{UDP}
      \small Sin acuses de recibo: satura el canal y cuantifica la pérdida de paquetes y el
      jitter frente a obstáculos o personas.
    \end{block}
  \end{columns}
  \vspace{4pt}
  \small Con PCs en los extremos, el tráfico lo generan las computadoras y las antenas
  funcionan solo como puente transparente, sin el cuello de botella de su CPU.
\end{frame}

\begin{frame}{Variación del ancho de canal: 60 MHz vs.\ 10 MHz}
  \begin{columns}[c]
    \column{0.42\textwidth}\centering
    \figura{panel_10mhz.jpg}{\linewidth}\\ {\scriptsize Panel a 10 MHz (alineadas)}
    \column{0.56\textwidth}\centering
    \footnotesize
    \begin{tabular}{@{}lcc@{}}
      \toprule
       & \textbf{60 MHz} & \textbf{10 MHz} \\ \midrule
      Señal (alineadas) & $-29$ dBm & $-26$ dBm \\
      Piso de ruido & $-86$ dBm & $-92$ dBm \\
      Capacidad & 454 Mbps & 64.7 Mbps \\
      Modulación (alin.) & 8x 256QAM & 8x 256QAM \\
      Modulación (desal.) & 2x QPSK & 6x 64QAM \\
      Ping mín.\ / máx. & 1.13 / 8.10 ms & 2.42 / 2.91 ms \\
      \bottomrule
    \end{tabular}
  \end{columns}
  \vspace{6pt}
  \small
  \begin{itemize}
    \item La señal no cambia (misma potencia y distancia); el ruido baja $\approx$6 dB
          ($10\log 6 = 7.8$ dB teóricos) y mejora el SNR.
    \item Un sexto del espectro $\Rightarrow$ capacidad $\approx$7 veces menor, aun con 256QAM.
    \item Canal angosto: más latencia base, pero mucho menos jitter frente al \emph{multipath}.
  \end{itemize}
\end{frame}

\begin{frame}{Obstrucción de la zona de Fresnel}
  \begin{columns}[c]
    \column{0.55\textwidth}
    \begin{itemize}
      \item No hizo falta introducir un obstáculo: el \textbf{tránsito de estudiantes}
            funcionó como obstrucción dinámica.
      \item A 5 GHz el cuerpo humano (alto contenido de agua) atenúa fuertemente la señal.
      \item Efectos observados:
      \begin{itemize}
        \item caídas de $-28$ a $-56$ dBm,
        \item fluctuaciones en la herramienta de alineación,
        \item picos de latencia (jitter) por retransmisiones.
      \end{itemize}
    \end{itemize}
    \column{0.42\textwidth}
    \begin{alertblock}{Relación con la Parte 1}
      Confirma por qué el diseño exige despejar al menos el \textbf{60\,\%} de la primera
      zona de Fresnel para sostener el MCS y la capacidad.
    \end{alertblock}
  \end{columns}
\end{frame}
